\documentclass[11pt,a4paper]{article}

\usepackage[utf8]{inputenc} % allow utf-8 input
\usepackage[T1]{fontenc}    % use 8-bit T1 fonts
\usepackage{hyperref}       % hyperlinks
\usepackage{url}            % simple URL typesetting
\usepackage{booktabs}       % professional-quality tables
\usepackage{amsfonts}       % blackboard math symbols
\usepackage{nicefrac}       % compact symbols for 1/2, etc.
\usepackage{microtype}      % microtypography
\usepackage{xcolor}         % colors
\usepackage{graphicx}       % graphics
\usepackage[normalem]{ulem} % strikethrough

\begin{document}

\section{East Asian list numbering}

Articles and clauses numbered with chineseCountingThousand, chineseCounting and decimalEnclosedCircle. Article numbering starts at 9 so the list crosses the 十 boundary.

\begin{enumerate}
\item 合同双方应当按照约定履行各自的义务。
  \begin{enumerate}
\item 甲方应当按期支付价款。
\item 乙方应当按期交付标的物。
    \begin{enumerate}
\item 交付地点为甲方指定的仓库。
\item 交付时间为合同生效后三十日内。
    \end{enumerate}
  \end{enumerate}
\item 任何一方不得擅自变更或者解除合同。
\item 本合同自双方签字盖章之日起生效。
\end{enumerate}

\subsection{Other formats}

chineseLegalSimplified:

\begin{enumerate}
\item 第一项内容
\item 第二项内容
\item 第三项内容
\end{enumerate}

ideographTraditional:

\begin{enumerate}
\item 甲项
\item 乙项
\item 丙项
\end{enumerate}

decimalFullWidth:

\begin{enumerate}
\item 第一行
\item 第二行
\end{enumerate}

\end{document}